\chapter{Analysis}
\label{chap:analysis}
\graphicspath{{Chapter4/img/}}
\setlength{\emergencystretch}{2em}

% This chapter interprets the causal NLP results reported in
% Section~\ref{sec:impl-causal-nlp}. Data preparation, synthetic validation,
% exploratory analysis, and topic modelling are not repeated as independent analytical
% sections because they do not establish the policy--sentiment causal gap. Their role
% was to define the evidence boundary, preserve provenance, construct comparable units,
% and provide the frozen NMF topic spaces used by the causal pipelines. They form the
% measurement foundation for the analysis rather than separate answers to the research
% question.

% The term causal gap refers here to differences in causal language and discourse
% organisation. It does not establish real-world causal effects or policy failure.
% The analysis asks whether policy and original sentiment express comparable causal
% claims and whether they connect related topics through similar causes, mechanisms,
% requirements, and outcomes.

% \section{Interpretation of the Policy--Sentiment Causal Gap}
% \label{sec:analysis-gap}

% \subsection{Semantic coverage}

% The balanced global deficits were 0.3585 for policy-to-sentiment coverage and 0.3675
% in the reverse direction. Their difference of approximately 0.0089 is too small to
% support a strong directional asymmetry. The more important result is that both
% corpora contain broad semantic overlap alongside a persistent residual gap: many
% causal claims have related counterparts, but the correspondence is incomplete in
% both directions.

% The topic-level results locate this gap more clearly. P8, covering AI curriculum
% development and school oversight, had the largest policy-topic deficit at 0.3960.
% P0, covering AI-assisted learning and assessment protocols, followed at 0.3763 and
% was supported by 604 claims from 29 documents. P4 reached 0.3729 but drew on only six
% documents, so it is a narrower lead. P3 had the lowest deficit at 0.3079, indicating
% the closest semantic correspondence. Among sentiment topics, the largest deficits
% occurred for online-learning clarity, GenAI assessment confusion, and public
% perceptions of AI and youth employment. The semantic gap is therefore concentrated
% around curriculum oversight, assessment, and selected experiences rather than being
% uniform across the corpora.

% \subsection{Causal-frame organisation}

% The frame network provides stronger evidence of organisational difference. The
% primary Jensen--Shannon divergence was 0.2596, while the source-balanced,
% high-quality-only, and directly projected-span specifications ranged from 0.2425 to
% 0.2964. The gap therefore remained visible when source contribution, frame quality,
% and projection reliability changed. Although approximately 31--35\% of span
% assignments used document-topic fallback, the directly projected specification
% produced the largest divergence, showing that fallback use alone does not explain
% the result.

% The relation shares clarify the nature of the difference. Causes/increases accounted
% for 36.8\% of sentiment frame weight but 24.7\% of policy weight. Policy placed more
% weight on enables/supports (41.9\% versus 37.2\%), expected improvement (12.9\%
% versus 6.3\%), and requirements/dependencies (9.4\% versus 6.8\%). Risk/threat
% shares were similar. Policy discourse is therefore organised more around
% institutional mechanisms, obligations, and intended outcomes, whereas sentiment
% discourse gives more weight to experienced causes and consequences. The distinction
% is not simply positive policy language versus negative opinion; it concerns the
% causal role assigned to related themes.

% At topic level, P0 had the largest structural gap at 0.3109, followed by P2 at
% 0.3099 and P8 at 0.2959. P3 again had the smallest gap at 0.1301. The recurrence of
% P0 and P8 across the semantic and frame pipelines provides the strongest convergent
% evidence: these topics differ both in claim-level coverage and in how their causes
% and effects are organised.

% \subsection{Cross-method interpretation}

% \begin{table}[!htbp]
% \centering
% \small
% \setlength{\tabcolsep}{4pt}
% \begin{tabular}{
% >{\centering\arraybackslash}p{0.07\linewidth}
% >{\raggedright\arraybackslash}p{0.25\linewidth}
% >{\raggedright\arraybackslash}p{0.25\linewidth}
% >{\raggedright\arraybackslash}p{0.33\linewidth}
% }
% \hline
% Topic & Semantic evidence & Frame evidence & Interpretation \\
% \hline
% P8 & Largest policy-topic deficit & Third-largest structural gap &
% Convergent gap in curriculum development and school oversight. \\
% P0 & Second-largest deficit with broad support & Largest structural gap &
% Strongest combined gap in learning and assessment protocols. \\
% P2 & Less prominent semantic gap & Second-largest structural gap &
% Related claims exist, but their causal organisation differs. \\
% P4 & High deficit with narrow support & Lower structural prominence &
% Focused semantic lead requiring source-level inspection. \\
% P3 & Lowest semantic deficit & Lowest structural gap &
% Strongest alignment across both causal measures. \\
% \hline
% \end{tabular}
% \caption{Cross-method interpretation of the principal policy-topic gaps.}
% \label{tab:analysis-causal-synthesis}
% \end{table}

% The two methods should not be collapsed into one score. Semantic coverage identifies
% claims with weak counterparts, while frame divergence identifies different causal
% organisation even when related claims are present. P0 and P8 are the strongest
% findings because both methods identify them across complementary levels and their
% evidence spans broad document sets. P2 is mainly a framing gap, while P4 is mainly a
% semantic gap with limited source support. P3 provides the clearest case of
% cross-method alignment.

% Country results remain secondary. France showed large differences around P2 and P8,
% with P2 reaching 1.6518, while Ireland was led by P7 at 0.5391. However, France and
% Ireland each contributed only three original sentiment sources to the frame
% comparison. These outputs identify where source-level investigation is most useful;
% they do not support national rankings.

% Synthetic evidence also remains diagnostic. The semantic check changed most topic
% deficits only slightly, but S2 changed by approximately \(-0.057\) and S8 had no
% synthetic assignment. The synthetic frame network increased divergence from 0.2596
% to 0.3611 and altered requirements/dependencies most strongly. Synthetic data
% therefore tests sensitivity but does not validate the original causal structure.

% Overall, the evidence supports a selective rather than universal policy--sentiment
% gap. Policy and sentiment often discuss related AI-in-education concerns, yet they
% differ most around curriculum oversight, learning and assessment, and the causal
% organisation of institutional action versus experienced consequences. The result is
% strongest where semantic and structural evidence converge, and weaker where source
% coverage or projection confidence is limited.

% \section{LLM-Assisted Causal-Gap Interpretation}
% \label{sec:analysis-llm-causal-gap}

% This section will interpret the LLM-assisted outputs reported in
% Section~\ref{sec:impl-llm-causal-gap} and relate them to the semantic and frame-based
% findings above.

% % Add the validated LLM-assisted findings here.
% % Distinguish model-generated interpretation from corpus-grounded evidence.
% % Retain claim, frame, topic, country, and source identifiers for every conclusion.
